\documentclass[a4paper]{article}
\usepackage{amsfonts}    
\usepackage{amsmath}
\usepackage{amssymb}
\usepackage{graphicx}

\begin{document}
  
\noindent \large Auteur: Michael Zielinski \\
\noindent \large Studierichting: Informatica \\
\noindent \large Oefeningenreeks 6E nr. 8.2 \\

\medskip

\normalsize

Bereken met een Taylorreeks tot op $12$ decimalen: \\

$ \cos\left(\dfrac{1}{3}\right) $ \\

De Taylorreeks van $\cos x$ rond $0$ is: \\

$ \cos x = \displaystyle\sum_{n=0}^{\infty}(-1)^n\dfrac{x^{2n}}{(2n)!} $ \\

Dus: \\

$ \cos x = 1-\dfrac{x^2}{2!}+\dfrac{x^4}{4!}-\dfrac{x^6}{6!}+\dfrac{x^8}{8!}-\dfrac{x^{10}}{10!}+\dfrac{x^{12}}{12!}-\cdots $ \\

Voor $x=\dfrac{1}{3}$ krijgen we: \\

$ \cos\left(\dfrac{1}{3}\right)
= 1-\dfrac{\left(\dfrac{1}{3}\right)^2}{2!}
+\dfrac{\left(\dfrac{1}{3}\right)^4}{4!}
-\dfrac{\left(\dfrac{1}{3}\right)^6}{6!}
+\dfrac{\left(\dfrac{1}{3}\right)^8}{8!}
-\dfrac{\left(\dfrac{1}{3}\right)^{10}}{10!}
+\cdots $ \\

We moeten weten tot hoeveel termen we moeten gaan. \\

Omdat dit een alternerende reeks is, is de fout kleiner dan de eerstvolgende term die we niet meenemen. \\

Na de term met macht $10$ is de eerstvolgende term: \\

$ \dfrac{\left(\dfrac{1}{3}\right)^{12}}{12!} $ \\

$ = \dfrac{1}{3^{12}\cdot 12!} $ \\

$ = \dfrac{1}{531441\cdot 479001600} $ \\

$ \approx 3.93\cdot 10^{-18} $ \\

Omdat: \\

$ 3.93\cdot 10^{-18} < 0.5\cdot 10^{-12} $ \\

volstaat het om te rekenen tot en met de term met macht $10$. \\

Dus: \\

$ \cos\left(\dfrac{1}{3}\right)
\approx 1-\dfrac{1}{18}
+\dfrac{1}{1944}
-\dfrac{1}{524880}
+\dfrac{1}{264539520}
-\dfrac{1}{214326028800} $ \\

$ \approx 0.944956946315 $ \\

Dus: \\

$ \boxed{\cos\left(\dfrac{1}{3}\right) \approx 0.944956946315} $ \\

\begin{thebibliography}{999}
%\bibitem{a} Referentie
\end{thebibliography}

\end{document}